\documentclass[11pt]{article}
\usepackage[margin=0.9in]{geometry}
\usepackage{amsmath,amssymb,amsthm,booktabs}
\usepackage[T1]{fontenc}
\usepackage{lmodern,microtype}
\usepackage[colorlinks=true,linkcolor=blue,urlcolor=blue]{hyperref}
\newtheorem{theorem}{Theorem}[section]
\newtheorem{lemma}[theorem]{Lemma}
\newtheorem{proposition}[theorem]{Proposition}
\newtheorem{corollary}[theorem]{Corollary}
\theoremstyle{remark}
\newtheorem{remark}[theorem]{Remark}
\title{Toward a four-dimensional Yang--Mills construction\\
\large Notebook 39: normalized local replacement in the bare Wilson measure}
\author{Research proof draft}
\date{17 October 2026}
\setlength{\emergencystretch}{5em}
\begin{document}
\maketitle
\begin{abstract}
We complete a normalized replacement operation for arbitrary finite
collections of separated fixed boxes in the positive bare Wilson
measure. On a good boundary the exact conditional law is restricted
to its local gauge/minimizer window and divided by its actual
conditional probability. On a bad boundary the law is unchanged.
The operation preserves the exterior marginal, gauge invariance,
and total raw partition mass. A positive expansion accounts for
every discarded interior configuration.

An explicit coupling has discrepancy marks dominated by independent
Bernoulli variables of parameter \(a=Ce^{-cT^2}\). For the exterior
jointly with any selected set \(A\) of box interiors, the exact
total variation distance is
\(\mathbb E[1-\prod_{i\in A}(1-p_i)]\le a|A|\).
We also prove local entropy, sparse-defect, cluster,
bounded-source and connected-cumulant comparison bounds. All are
independent of the number of other boxes.

The probability inputs are the full normalized Wilson bounds of
Notebook 38, not Gaussian surrogates. Geometry and box size remain
fixed. The operation has not been identified with a nonlinear
coarse-average fiber or propagated through general mixed
renormalization histories. This completes a restricted
finite-measure theorem and gives a definite stopping point for
handoff; it does not complete continuum Yang--Mills existence.
\end{abstract}

\section*{Inputs and proof boundary}

The two quantitative inputs are the fixed-box conditional tail
theorem and joint boundary-plaquette probability theorem
proved in Notebook 38. Both concern the actual compact positive
Wilson measure and include their normalization. Their box
geometry, boundary gauge, adjoint-invariant window, period
conditions and constants are restated below.
The new argument is an exact normalized operation built from
those inputs, together with its local comparison estimates.
It does not use the earlier homogeneous determinant calculation.

The group may be \(SU(N)\) for these bare-measure statements,
with \(\tau=\operatorname{Tr}/N\). This does not extend the
earlier \(SU(2)\) multiscale construction to every group.
An internal audit is included in the research archive.
Such review and successful typesetting do not certify
all upstream arguments independently.



% No preamble. Prefix nbr-. 17 October 2026.
\section{A normalized fixed-box replacement for the full Wilson measure}
\label{nbr-section}

This section proves a restricted normalized replacement theorem for the
\emph{actual positive bare Wilson measure}. It preserves the exterior
marginal and the raw partition mass exactly, leaves bad-boundary boxes
unchanged, and gives a sparse failure expansion and local comparison
bounds uniform in the number of boxes. The fixed-box coordinates are
not identified with a nonlinear coarse-average fiber.

\subsection{Exact conditional structure and the established box estimates}
\label{nbr-input-section}

Let
\[
 \mu(dU)=Z_W^{-1}e^{-S_W(U)/g^2}\prod_e dU_e
\]
be the normalized positive Wilson measure on a finite compatible
four-dimensional periodic lattice, with normalized link Haar measure.
Fix \(R\ge2\), \(b\ge3\), and a finite collection \(I\) of embedded
translates \(\Omega_i\) of \([0,R]^4\), with pairwise disjoint vertex
sets and no local periodic identifications inside a box. The surrounding
periods are divisible by \(2b\). Constants may depend on this fixed
geometry and group, but not on the number or placement of boxes or on
the surrounding volume.

Let \(E_{\circ,i}\) be the edges of box \(i\) not contained in its
topological boundary. The exterior variable \(\omega\) comprises every
edge outside their union, including all box boundary edges; let
\(\mu_{\rm ext}\) be its actual marginal. Write \(u_i\) for the
interior edges. Every plaquette meeting \(E_{\circ,i}\) lies entirely
in \(\Omega_i\). No plaquette therefore depends on two integrated
interiors. The exact conditional law is
\begin{equation}
 \mu(d\omega,du)=\mu_{\rm ext}(d\omega)
                    \prod_{i\in I}Q_i^\omega(du_i),
 \label{nbr-disintegration}
\end{equation}
where \(Q_i^\omega\) is the normalized compact Wilson law of that box
with its actual boundary fixed. It depends only on that boundary.
All exterior action terms cancel from this normalization.

Use the boundary-only tree gauge and then the forest rooted at the
boundary vertices from the fixed-box result. Let \(B_i(\omega)\) be
the event that some boundary plaquette has
\(h(U_p)\ge\varepsilon_R^2\). On \(B_i^c\), the transformed boundary
lies in the fixed small-boundary neighborhood. Let \(W_i(\omega)\)
be the lifted window about its unique global conditional minimizer,
\begin{equation}
 \|x-x_*(V_i(\omega))\|_\infty\le gT,
 \label{nbr-window}
\end{equation}
using the invariant edge/color norm in the forest coordinates.
It is an actual measurable gauge-invariant event: the full tree
representative and its unique minimizer transform by the same constant
root conjugation.

The already proved fixed-box conditional theorem supplies constants
\(C,c,r,g_0>0\) such that
\begin{equation}
 Q_i^\omega(W_i^c)\le a:=Ce^{-cT^2}\quad\hbox{on }B_i^c,
 \qquad 0<g\le g_0,\quad 1\le T\le r/g.
 \label{nbr-conditional-tail}
\end{equation}
Choose a fixed \(T_0\ge1\) so that \(Ce^{-cT_0^2}\le1/2\), and
henceforth take \(T_0\le T\le r/g\). Thus \(a\le1/2\).
This estimate uses the entire compact conditional integral: its
\(g^n e^{-S_{\min}/g^2}\) denominator and local numerator powers
cancel, and its complementary compact region has a positive action
gap. It is not a Gaussian-law substitution.

The full Wilson boundary estimate supplies, for every \(S\subset I\),
\begin{equation}
 \mu_{\rm ext}\!\left(\bigcap_{i\in S}B_i\right)
 \le\beta^{|S|},\qquad
 \beta=\min\{1,24R^3z_b(g,\varepsilon_R)^{1/Q_b}\},
 \quad Q_b=6b^4,
 \label{nbr-boundary-tail}
\end{equation}
where
\[
 z_b(g,\delta)=\min\left\{1,
 \exp\left[-\delta^2/g^2+
 b^4\{4\dim G\log(1/g)+C_{G,4}\}\right]\right\}.
\]
Indeed the boundary events are exterior measurable. A union bound over
one offending boundary plaquette in each box and the prescribed-
plaquette Wilson estimate prove \eqref{nbr-boundary-tail}; the chosen
plaquettes are distinct by vertex disjointness. Boundary independence
is not assumed. These are fixed-geometry inputs from the full-measure
box theorem, not assertions for arbitrary effective actions.

\subsection{The normalized operation}
\label{nbr-operation-section}

Define
\begin{equation}
 p_i(\omega)=\mathbf1_{B_i^c}Q_i^\omega(W_i^c),\qquad
 P_i^\omega=
 \begin{cases}
 Q_i^\omega(\,\cdot\mid W_i),&B_i^c,\\
 Q_i^\omega,&B_i.
 \end{cases}
 \label{nbr-kernels}
\end{equation}
On good boundaries, \(1-p_i\ge1-a\ge1/2\); on bad boundaries
\(p_i=0\) and the conditional law is left unchanged. All kernels are
measurable finite-dimensional compact integrals. For \(J\subset I\),
put
\begin{equation}
 \nu_J(d\omega,du)=\mu_{\rm ext}(d\omega)
       \prod_{i\in J}P_i^\omega(du_i)
       \prod_{i\notin J}Q_i^\omega(du_i),\qquad \nu=\nu_I.
 \label{nbr-replacement}
\end{equation}

\begin{theorem}[Normalized fixed-box replacement]
\label{nbr-main}
The measures \(\nu_J\) are probability measures with exactly the
original exterior marginal and with the Radon--Nikodym derivative
\begin{equation}
 \frac{d\nu_J}{d\mu}
 =\prod_{\substack{i\in J\\B_i^c}}
             \frac{\mathbf1_{W_i}(u_i)}{1-p_i(\omega)}.
 \label{nbr-rn}
\end{equation}
They are gauge invariant. The same replacement applied to the
unnormalized Wilson measure preserves its total mass \(Z_W\) exactly.
There is a coupling of \(\mu\) and \(\nu\) with the same exterior
whose box discrepancy marks are dominated by independent Bernoulli
variables of parameter \(a\). Bad-boundary interiors are identical
in this coupling. In particular every fixed local observable admits
comparison bounds independent of \(|I|\), as quantified below.
\end{theorem}
\begin{proof}
Conditioning on \(\omega\) and multiplying the normalized kernels
proves \eqref{nbr-rn}. Each factor has conditional \(Q_i^\omega\)
expectation one. Thus the product has expectation one and
\begin{equation}
 \int e^{-S_W/g^2}
       \prod_{\substack{i\in J\\B_i^c}}
          \frac{\mathbf1_{W_i}}{1-p_i}\prod_e dU_e=Z_W.
 \label{nbr-raw-mass}
\end{equation}
Every exterior observable is preserved exactly. Gauge covariance of
conditional Haar/Wilson laws, invariance of \(B_i\), and invariance
of the lifted tree-gauge window show that \eqref{nbr-rn} is gauge
invariant. The explicit coupling is constructed in the next subsection.
\end{proof}

The conditional denominators are exact boundary-dependent quantities,
not Gaussian determinants or vacuum constants. They obey
\begin{equation}
 0\le-\log(1-p_i)\le\frac{p_i}{1-a}\le\frac a{1-a}.
 \label{nbr-denominator-bound}
\end{equation}
Their dependence is confined to the fixed box boundary. The full
product density may nevertheless grow exponentially with \(|J|\);
no uniform global density-ratio bound is asserted.
The individual kernels also define resampling operators: keep the
exterior and redraw interior \(i\) from \(P_i^\omega\). These
operators commute for distinct boxes and are idempotent, because no
resampling changes another box's boundary. Their product applied once
to \(\mu\) gives \(\nu\).

\subsection{An exact positive failure expansion and coupling}
\label{nbr-coupling-section}

On a good boundary with \(p_i>0\), let
\(F_i^\omega=Q_i^\omega(\,\cdot\mid W_i^c)\).
If \(p_i=0\), choose any probability kernel for \(F_i\); it will be
multiplied by zero. For every \(\omega\),
\begin{equation}
 Q_i^\omega=(1-p_i)P_i^\omega+p_iF_i^\omega.
 \label{nbr-one-mixture}
\end{equation}
Consequently
\begin{align}
 \mu&=\sum_{S\subset I}\mathcal M_S,\notag\\
 \mathcal M_S(d\omega,du)
 &=\mu_{\rm ext}(d\omega)
   \prod_{i\in S}p_i\prod_{i\notin S}(1-p_i)
   \prod_{i\in S}F_i^\omega(du_i)
   \prod_{i\notin S}P_i^\omega(du_i).
 \label{nbr-failure-expansion}
\end{align}
All terms are positive and their masses sum to one. The empty term is
\(\prod_i(1-p_i(\omega))\,\nu\), not in general a constant multiple
of \(\nu\). Thus no exterior reweighting is hidden in this expansion.

To couple the measures, sample \(\omega\sim\mu_{\rm ext}\).
Conditionally and independently across boxes, sample
\(Y_i\sim P_i^\omega\), an independent Bernoulli mark \(D_i\) with
parameter \(p_i\), and, when needed, \(Z_i\sim F_i^\omega\). Set
\begin{equation}
 X_i=\begin{cases}Y_i,&D_i=0,\\Z_i,&D_i=1.\end{cases}
 \label{nbr-coupling}
\end{equation}
Then \((\omega,X)\sim\mu\) and \((\omega,Y)\sim\nu\).
On a marked box, \(X_i\notin W_i\) and \(Y_i\in W_i\); elsewhere
they are identical. Therefore \(D_i\) is precisely its discrepancy
indicator. On bad boundaries \(p_i=0\), so the interior remains the
same sample from the original law.

Take independent uniform random variables \(U_i\), independent of the
exterior, and realize \(D_i=\mathbf1_{U_i\le p_i(\omega)}\).
Then \(D_i\le\mathbf1_{U_i\le a}\) simultaneously. This proves the
iid Bernoulli domination asserted in Theorem~\ref{nbr-main}, and
\begin{equation}
 \mathbb P(D_i=1\text{ for all }i\in S)
 =\mathbb E\prod_{i\in S}p_i\le a^{|S|}.
 \label{nbr-discrepancy-joint}
\end{equation}
This stochastic domination concerns the coupling marks, not the
possibly correlated bad boundaries.

\subsection{Local total variation, observables and entropy}
\label{nbr-local-section}

Use the convention
\(\|\lambda-\lambda'\|_{\rm TV}=\sup_E|\lambda(E)-\lambda'(E)|\).
For \(A\subset I\), retain the whole exterior but only the interiors
in \(A\). Conditionally, their replacement law is their original law
conditioned on all its good-boundary windows, of probability
\(\prod_{i\in A}(1-p_i)\). Hence
\begin{equation}
 \|\mu_{({\rm ext},A)}-\nu_{({\rm ext},A)}\|_{\rm TV}
 =\mathbb E_{\mu_{\rm ext}}
       \left[1-\prod_{i\in A}(1-p_i)\right]
 \le a|A|.
 \label{nbr-local-tv}
\end{equation}
The equality follows by integrating the conditional positive and
negative density parts; equivalently, the coupling is identical with
that conditional probability, and any discrepancy puts the original
field outside the replacement support. For the marginal on \(A\)
alone the same expression is an upper bound.

For a bounded observable supported on the exterior and these interiors,
\begin{equation}
 |\mu(f)-\nu(f)|\le2\|f\|_\infty a|A|.
 \label{nbr-local-observable}
\end{equation}
More generally, let \(\operatorname{osc}_i(f)\) be the supremum of the
change when only box \(i\)'s interior changes. Telescoping along the
coupling proves
\begin{equation}
 |\mu(f)-\nu(f)|
 \le\sum_i\operatorname{osc}_i(f)\,\mathbb E p_i
 \le a\sum_i\operatorname{osc}_i(f).
 \label{nbr-oscillation}
\end{equation}
Thus local error depends on the region read by the observable, not on
the total number of replaced boxes. Global total variation can approach
one as that number grows.

The local joint-with-exterior relative entropy is also exact:
\begin{equation}
 D(\nu_{({\rm ext},A)}\|\mu_{({\rm ext},A)})
 =\mathbb E_{\mu_{\rm ext}}\sum_{i\in A}-\log(1-p_i)
 \le\frac{|A|a}{1-a}.
 \label{nbr-entropy}
\end{equation}
Substitute the local version of \eqref{nbr-rn}; all its window
indicators equal one under \(\nu\). No assertion about finiteness of
the reverse entropy is needed.

\subsection{Untreated boundaries, exponential moments and clusters}
\label{nbr-sparsity-section}

The exterior marginal is unchanged, so
\begin{equation}
 \nu\!\left(\bigcap_{i\in S}B_i\right)
 =\mu\!\left(\bigcap_{i\in S}B_i\right)\le\beta^{|S|}.
 \label{nbr-retained-boundaries}
\end{equation}
Under \(\nu\), the original box defect (bad boundary or a failed window
on a good boundary) is exactly \(B_i\) almost surely. Its good-boundary
part has been removed by an exactly normalized conditional operation;
its bad-boundary law has not been discarded or altered.

On the coupling put \(C_i=\mathbf1_{B_i}\lor D_i\). These alternatives
are disjoint at one box. Expanding a prescribed intersection over the
subset assigned to the boundary alternative, and conditioning on the
exterior, gives
\begin{align}
 \mathbb P(C_i=1\text{ for all }i\in S)
 &\le\sum_{J\subset S}a^{|S|-|J|}
       \mu_{\rm ext}\!\left(\bigcap_{i\in J}B_i\right)
 \le(a+\beta)^{|S|}.
 \label{nbr-combined-joint}
\end{align}
Let \(\varepsilon=\min(1,a+\beta)\).
Expansion of the finite product proves, for \(t\ge0\),
\begin{equation}
 \mathbb E\exp\left(t\sum_{i\in A}C_i\right)
 \le[1+\varepsilon(e^t-1)]^{|A|}.
 \label{nbr-mgf}
\end{equation}
The discrepancy version has \(a\) in place of \(\varepsilon\), even
conditionally on the exterior. This joint-inclusion argument does not
assert iid stochastic domination for \(C\).
For \(1\le k\le|A|\), a union bound gives
\begin{equation}
 \mathbb P\!\left(\sum_{i\in A}C_i\ge k\right)
 \le\min\left\{1,\binom{|A|}{k}\varepsilon^k\right\}
 \le\min\{1,(e|A|\varepsilon/k)^k\}.
 \label{nbr-count-tail}
\end{equation}
For \(k>|A|\) the probability is zero. One may also optimize the
Chernoff estimate from \eqref{nbr-mgf}.

On any fixed adjacency graph of boxes with maximum degree
\(\Delta\ge2\), there are at most \(\Delta^{2m}\) connected
\(m\)-subsets containing a given vertex: encode a canonical spanning
tree by its depth-first traversal. A cluster of size at least \(m\)
contains such a subset. Therefore
\begin{equation}
 \mathbb P(|C\text{-cluster at }i|\ge m)
       \le[\Delta^2\varepsilon]^m.
 \label{nbr-cluster}
\end{equation}
The same argument gives \(a\) for discrepancies and \(\beta\) for
defects under the replacement measure. These estimates are uniform
in volume and the total number of boxes, at the fixed geometry/degree.
For \(T(g)=(\log(1/g))^\alpha\), \(\alpha>1/2\), both \(a\)
and \(\beta\) are smaller than every fixed power of \(g\) at small
coupling. All fixed-region bounds inherit that property.

\subsection{Exact source identities and local connected cumulants}
\label{nbr-source-section}

For bounded measurable single-box sources \(f_i(\omega,u_i)\) and an
exterior source \(f_{\rm ext}(\omega)\), conditional independence gives
\begin{align}
 \nu\!\left(e^{f_{\rm ext}+\sum_i f_i}\right)
 &=\int\mu_{\rm ext}(d\omega)e^{f_{\rm ext}}
                             \prod_iP_i^\omega(e^{f_i}),
 \label{nbr-source-replacement}\\
 \mu\!\left(e^{f_{\rm ext}+\sum_i f_i}\right)
 &=\int\mu_{\rm ext}(d\omega)e^{f_{\rm ext}}
       \prod_i\bigl[(1-p_i)P_i^\omega(e^{f_i})
                           +p_iF_i^\omega(e^{f_i})\bigr].
 \label{nbr-source-original}
\end{align}
On good boundaries,
\(P_i^\omega(e^{f_i})=Q_i^\omega(\mathbf1_{W_i}e^{f_i})/(1-p_i)\);
on bad boundaries it is unchanged. The exact zero-source conditional
masses remain in these denominators. No exterior-dependent factor is
absorbed into a vacuum constant. The identities are positive expansions
for real sources and absolutely valid for bounded complex sources.
Cross-box sources use the full measure expansion, not a false product
factorization in \eqref{nbr-source-replacement}.

For any bounded \(F\) depending on the exterior and interiors in \(A\),
\(\|F\|_\infty\le M\), the coupling yields
\begin{equation}
 |\mu(e^F)-\nu(e^F)|\le2e^M a|A|.
 \label{nbr-source-difference}
\end{equation}
For real \(F\), both expectations are at least \(e^{-M}\), so their
real logarithms differ by at most \(2e^{2M}a|A|\).

Now let \(O_1,\ldots,O_q\) be bounded observables with the same fixed
local dependence, put \(M_l=\|O_l\|_\infty>0\), and omit zero
observables. Fix
\(0<s_0<\log2\), \(d_0=2-e^{s_0}>0\).
For \(F(t)=\sum_lt_lO_l\) and
\(Z_\lambda(t)=\lambda(e^{F(t)})\), the weighted source ball
\(\sum_l|t_l|M_l\le s_0\) satisfies
\begin{equation}
 |Z_\mu(t)-1|,\ |Z_\nu(t)-1|\le e^{s_0}-1<1.
 \label{nbr-source-disk}
\end{equation}
Both functions are nonzero, with analytic logarithms zero at the
origin on the same branch. The straight segment between their values
also remains in this disk. Integrating \(1/z\) along it and using
\eqref{nbr-source-difference} gives
\begin{equation}
 |\log Z_\mu(t)-\log Z_\nu(t)|
       \le\frac{2e^{s_0}}{d_0}\,a|A|.
 \label{nbr-log-source}
\end{equation}
For another bounded observable \(O\) whose support is included in
\(A\), the normalized source numerator obeys
\begin{equation}
 \left|\frac{\mu(Oe^F)}{Z_\mu}
             -\frac{\nu(Oe^F)}{Z_\nu}\right|
 \le\frac{4\|O\|_\infty e^{s_0}}{d_0^2}\,a|A|.
 \label{nbr-source-ratio}
\end{equation}
Indeed the numerator difference is at most
\(2\|O\|_\infty e^{s_0}a|A|\), the denominator difference at most
\(2e^{s_0}a|A|\), their moduli are at least \(d_0\), and
\(|\nu(Oe^F)|\le\|O\|_\infty e^{s_0}\). Use
\(d_0+e^{s_0}=2\) to obtain the displayed constant.

Let \(k=(k_1,\ldots,k_q)\ne0\), \(n=\sum_l k_l\).
Cauchy on radii \(s_0/(qM_l)\) gives the explicit connected-cumulant
comparison
\begin{equation}
 \left|\partial^k(\log Z_\mu-\log Z_\nu)(0)\right|
 \le\frac{2e^{s_0}}{d_0}\,a|A|\,
           k!\left(\frac q{s_0}\right)^n\prod_l M_l^{k_l}.
 \label{nbr-cumulant}
\end{equation}
Alternatively set inactive variables to zero and use radii
\(s_0k_l/(nM_l)\). The factor preceding \(\prod M_l^{k_l}\)
then improves from \(k!(q/s_0)^n\) to
\[
 \frac{k!(n/s_0)^n}{\prod_{k_l>0}k_l^{k_l}}
       \le n!(e/s_0)^n.
\]
The last inequality uses \(k!\le\prod k_l^{k_l}\) and
\(n^n\le e^n n!\). All finite generating functions are entire for
bounded observables, so differentiation under their integrals is
valid. Windows and coupling are held fixed during differentiation.
These statements are for bounded observables with fixed local support;
uniformity in a later cutoff applies only if their displayed bounds
remain uniform. They do not constitute estimates for renormalized
unbounded fields.

\subsection{Natural stopping point and scope}
\label{nbr-scope}

This operation treats an arbitrary finite union of vertex-disjoint
fixed boxes in the full positive bare measure. Its conditional
normalizations are exact; it preserves the exterior marginal, total
partition mass and gauge invariance; and it accounts for every removed
configuration through a positive sparse expansion. Unlike conditioning
the whole lattice on a simultaneous good event, it does not reweight
the exterior distribution. It also supplies local comparison, source,
cumulant and cluster bounds uniform in the number of boxes.

The restriction is substantive. The box size and forest-coordinate
conditional minimizer are fixed. No identification with an averaged
coarse-constraint fiber, arbitrary growing component, collared
multiscale determining set, scale-age estimate or effective action
after renormalization is claimed. Bad-boundary boxes remain with their
original conditional laws. Complex source formulas here come from
bounded finite observables, not a complex-measure stability theorem.
The result is a normalized local replacement input for future
mixed-history analysis; it does not close the continuum existence,
reconstruction, nontriviality or mass-gap problem.

\end{document}
